\documentclass[11pt,a4paper]{article}

\usepackage[margin=24mm]{geometry}
\usepackage{fontspec}
\usepackage{xeCJK}
\usepackage{microtype}
\usepackage{booktabs}
\usepackage{longtable}
\usepackage{tabularx}
\usepackage{array}
\usepackage{xcolor}
\usepackage{enumitem}
\usepackage{hyperref}
\usepackage{url}
\usepackage{fancyhdr}
\usepackage{titlesec}
\usepackage{listings}

\setmainfont{Times New Roman}
\setsansfont{Arial}
\setmonofont{Menlo}
\setCJKmainfont{Songti SC}
\setCJKsansfont{PingFang SC}
\setCJKmonofont{PingFang SC}

\hypersetup{
  colorlinks=true,
  linkcolor=black,
  urlcolor=blue,
  citecolor=black,
  pdfauthor={SolvingLab},
  pdftitle={OmniWeave 与代码图谱/代码库记忆生态深度调研},
  pdfsubject={codegraph, codebase memory, MCP, code intelligence},
  pdfkeywords={OmniWeave, codegraph, MCP, code intelligence, codebase memory}
}

\definecolor{owblue}{HTML}{1F4E79}
\definecolor{owgray}{HTML}{F4F5F7}
\definecolor{owline}{HTML}{D0D7DE}
\definecolor{owgreen}{HTML}{1F7A4D}
\definecolor{owred}{HTML}{8A1F11}

\pagestyle{fancy}
\fancyhf{}
\lhead{\sffamily OmniWeave 调研}
\rhead{\sffamily \thepage}
\renewcommand{\headrulewidth}{0.3pt}

\titleformat{\section}
  {\Large\sffamily\bfseries\color{owblue}}
  {\thesection}{0.75em}{}
\titleformat{\subsection}
  {\large\sffamily\bfseries}
  {\thesubsection}{0.75em}{}
\titleformat{\subsubsection}
  {\normalsize\sffamily\bfseries}
  {\thesubsubsection}{0.75em}{}

\setlist[itemize]{leftmargin=1.6em, itemsep=0.25em, topsep=0.25em}
\setlist[enumerate]{leftmargin=1.8em, itemsep=0.25em, topsep=0.25em}
\renewcommand{\arraystretch}{1.2}
\setlength{\headheight}{14pt}
\emergencystretch=3em
\sloppy
\newcommand{\code}[1]{\texttt{\detokenize{#1}}}
\newcommand{\localpath}[1]{\path{#1}}
\newcommand{\evidence}[1]{\textcolor{owgreen}{#1}}
\newcommand{\risk}[1]{\textcolor{owred}{#1}}
\newcommand{\project}[1]{\textbf{#1}}

\lstset{
  basicstyle=\small\ttfamily,
  backgroundcolor=\color{owgray},
  frame=single,
  rulecolor=\color{owline},
  breaklines=true,
  columns=fullflexible,
  keepspaces=true,
  showstringspaces=false
}

\begin{document}

\begin{titlepage}
  \centering
  \vspace*{18mm}
  {\Huge\sffamily\bfseries OmniWeave 与代码图谱/代码库记忆生态深度调研\par}
  \vspace{8mm}
  {\Large\sffamily 我们真正的优势、边界，以及可借鉴路线\par}
  \vspace{16mm}
  {\large 研究日期：2026-06-23\par}
  \vspace{3mm}
  {\large 研究目录：\localpath{/Users/liuzaoqu/Desktop/develop/sogen/OmniWeave/research/2026-06-23-codegraph-ecosystem}\par}
  \vfill
  \begin{tabular}{rl}
    证据方式 & 源码 clone、关键文件逐行阅读、官方文档与论文交叉验证 \\
    研究对象 & 26 个项目/系统，18 个本地源码快照，8 类官方/论文资料 \\
    输出目标 & 帮 OmniWeave 做产品与架构决策，而不是项目百科 \\
  \end{tabular}
  \vfill
  {\small 本报告所有性能/覆盖数字均按来源分级处理：本地实测、项目自述、论文声明、官方文档声明严格分开；未复现实测的外部指标不作为 OmniWeave 决策事实。}
\end{titlepage}

\tableofcontents
\newpage

\section{结论先行}

\subsection{一句话结论}

OmniWeave 的核心优势不在“支持更多语言”、不在“有一个图数据库”、也不在“比 grep 更聪明”。真正的优势是：

\begin{quote}
  \textbf{它把 coding agent 最容易付出高成本、且 LSP/grep 结构性看不见的跨边界关系，做成了可追踪、可置信、可预算控制的本地图。}
\end{quote}

这类项目的生态正在分成三条路：第一条是 \project{结构图/代码智能后端}，代表是 codegraph、Codebase-Memory、CGC、Codanna、code-review-graph；第二条是 \project{语义记忆/向量检索}，代表是 Claude Context、semantic-search-mcp、Cursor 类索引；第三条是 \project{工业级精确代码索引/安全分析}，代表是 SCIP/Sourcegraph、Kythe、Glean、CodeQL、Joern、Semgrep/OpenGrep。OmniWeave 应该站在第一条，但吸收第二条作为入口、第三条作为可选事实来源，而不能被任何一条吞掉。

\subsection{OmniWeave 当前最强的五个点}

\begin{enumerate}
  \item \textbf{边界关系是第一公民。} OmniWeave 明确抓 Python/JS/TS/Go/workflow 到本地脚本、R S4 dispatch、Snakemake/Nextflow artifact DAG、external tool invocation。这不是“多语言列表”，而是 LSP 和 grep 的盲区建模。README 将这一点写得非常清楚：cross-language、cross-process、dynamic dispatch、workflow DAG、external-tool graph 是产品主线，见 \localpath{README.md:22-31,135-164}。
  \item \textbf{诚实边界强。} \code{generalCrossLangEdges} 对动态路径、插值 basename、非真实 subprocess 调用做 gate；推断边带 \code{provenance}、\code{metadata.confidence}、\code{registeredAt}，静态不可解就跳过，见 \localpath{src/resolution/callback-synthesizer.ts:1941-2039}。这比“什么都连上”的图更适合 agent，因为错边比漏边更危险。
  \item \textbf{价值经过 agent A/B 约束。} OmniWeave 的 README 和 eval 资料没有把工具包装成“更正确”，而是承认正确性常常与 grep/LSP 打平，优势主要在 reverse/blast-radius 的工具调用、token、延迟和稳定性，见 \localpath{README.md:35-98} 和 \localpath{eval-results/agent-ab-2026-06-13/}。
  \item \textbf{架构复杂度克制。} 当前主链是 WASM tree-sitter extraction → SQLite/FTS → resolution/synthesizer → MCP。没有默认外部向量库、没有多后端图数据库、没有把 LLM 文档生成塞进核心。CHECKPOINT 的模块边界和 skip-over-guess 原则与此一致，见 \localpath{CHECKPOINT.md:38-55,77-88}。
  \item \textbf{agent 输出面被当作产品，而不是 API 附属物。} 继承自 codegraph 的输出预算、\code{explore}/\code{node} 作为 primary tool、工具 surface 收缩等思想，使图不是“给人查的数据库”，而是“给 agent 消费的行动上下文”。codegraph 的 \code{src/mcp/tools.ts} 已经把默认工具压到 \code{explore}，OmniWeave README 当前默认五个核心工具，方向一致。
\end{enumerate}

\subsection{最大的战略边界}

OmniWeave 不应试图成为所有东西：

\begin{itemize}
  \item 不应替代 LSP 在同语言精确导航上的优势。LSP/SCIP 有编译器语义，OmniWeave 的价值在跨边界和零配置。
  \item 不应用向量检索伪装成结构事实。语义搜索适合“auth 在哪儿”这种概念入口，但不能直接产生 \code{calls}/\code{crossLang}/\code{overrides} 这类可置信边。
  \item 不应追求 Codebase-Memory/CGC 式大工具面。14 个、20 个、甚至更多 MCP tools 会提高试错概率；OmniWeave 更应该把少数工具做到足够完整。
  \item 不应为了“企业感”引入多图数据库、多 UI、LLM 文档生成、自然语言 Cypher 等核心依赖。那些是外围产品层，不是最小可信 substrate。
\end{itemize}

\section{研究方法与证据等级}

\subsection{本次调研范围}

本次建立了独立研究目录：

\begin{center}
\localpath{/Users/liuzaoqu/Desktop/develop/sogen/OmniWeave/research/2026-06-23-codegraph-ecosystem}
\end{center}

其中 \localpath{metadata/repos.tsv} 记录 26 个项目，\localpath{metadata/repo-snapshot.tsv} 记录 18 个本地源码快照的 commit、文件数、体量；\localpath{project-inventory.md} 记录证据等级。源码 clone 不是为了“多看几个 README”，而是为了验证每类系统真正的架构重心。

\subsection{证据等级}

\begin{longtable}{p{0.24\linewidth}p{0.68\linewidth}}
\toprule
等级 & 含义 \\
\midrule
\endhead
\code{source-clone} & 本地 clone 并阅读关键实现。本文可讨论源码结构、工具面、存储、索引管线、输出协议。 \\
\code{paper+source-clone} & 同时具备论文/项目叙述和可读源码；本文区分论文目标与实现质量。 \\
\code{source-sampled} & 没有完整 clone 或只取官方/代表性文件；仅讨论公开设计方向，不做实现细节断言。 \\
\code{official-docs} & 官方文档/博客/产品说明；可用于定位与能力边界，但数字按“官方声明”处理。 \\
\code{paper} & 论文指标或 benchmark；除非本地复现，否则不作为 OmniWeave 事实，只作为方向信号。 \\
\bottomrule
\end{longtable}

\subsection{反营销处理原则}

这类项目普遍喜欢宣称“百万行秒级索引”“数百倍 token 节省”“支持上百语言”。本报告采用三条过滤规则：

\begin{enumerate}
  \item \textbf{项目自述数字只标注来源，不直接变成结论。} 例如 Codebase-Memory 的 Linux kernel 级索引数字、Codebase-Memory paper 的 83\% answer quality、code-review-graph 的 82x/528x token reduction，都按来源分级。
  \item \textbf{语言数量不是主要指标。} 关键是该语言是否有可用关系、是否能跨文件/跨进程、是否有 false-positive gate、是否能被 agent 直接消费。
  \item \textbf{工具数量不是优势。} MCP 工具越多，agent 选择成本越高。真正优势是少数工具是否能一次返回可行动上下文。
\end{enumerate}

\section{生态分层：这类项目到底在竞争什么}

\subsection{四个真正的问题}

所有“code graph / codebase memory / code intelligence for agents”项目，本质上都在回答四个问题：

\begin{enumerate}
  \item \textbf{如何把代码变成可查询的事实？} 选择 Tree-sitter、LSP、SCIP、compiler indexer、regex、embedding，决定了事实的可信度。
  \item \textbf{如何把事实保持新鲜？} 选择 file watcher、git diff、hash/Merkle、incremental rebuild、snapshot artifact，决定了开发时可用性。
  \item \textbf{如何把事实交给 agent？} 选择 MCP tools、skills、slash commands、CLI、IDE extension、query language，决定了 agent 会不会用错。
  \item \textbf{如何证明它值得存在？} 选择 unit eval、agent A/B、retrieval benchmark、security finding、token accounting，决定了项目是否会陷入自我宣传。
\end{enumerate}

\subsection{生态分类表}

\begin{longtable}{p{0.18\linewidth}p{0.22\linewidth}p{0.24\linewidth}p{0.26\linewidth}}
\toprule
类型 & 代表项目 & 真正强项 & 对 OmniWeave 的意义 \\
\midrule
\endhead
Agent-first 结构图 & codegraph、OmniWeave、code-review-graph、Codanna & 少数查询直接给 agent 行动上下文；强调 token/tool-call 经济性 & 主战场。继续压低形态税、提高输出充分性。 \\
大型 codebase memory & Codebase-Memory、CGC & 多语言、大工具面、持久 DB、artifact、query language、安装面 & 借鉴 artifact/安全查询/控制面；避免工具面膨胀。 \\
语义检索/向量记忆 & Claude Context、semantic-search-mcp、Cursor indexing & 概念召回、自然语言入口、跨大代码库模糊搜索 & 做可选 sidecar：只产生候选 seed，不产生结构真边。 \\
LSP/IDE 能力层 & Serena、Blarify LSP path & find references、rename、diagnostics、precise edits & 做可选 validation/edit companion，不替代 core graph。 \\
工业精确索引 & SCIP/Sourcegraph、Kythe、Glean & 编译器级导航、跨 repo symbol、事实格式、Datalog/Angle 查询 & 借鉴事实格式和 importer；不把核心变成重型索引平台。 \\
安全/程序分析图 & CodeQL、Joern、Semgrep/OpenGrep & variant analysis、CPG、taint/rule matching & 可做未来安全插件，不是 coding-agent 日常上下文核心。 \\
研究型 RAG benchmark & RepoGraph、CodeRAG-Bench、CodexGraph/CodeRAG & 衡量检索对生成/修 bug 的影响 & 借鉴评测问题设计，不照搬实现。 \\
可视化学习工具 & Graphify、Understand Anything & onboarding、dashboard、文档/知识图、团队理解 & 可做外围报告/学习层；不要污染 MCP 核心。 \\
\bottomrule
\end{longtable}

\section{逐项目源码深读}

\subsection{codegraph：OmniWeave 的直接基座，真正价值在 agent 输出面}

\subsubsection{定位}

\project{codegraph} 不是最炫的图数据库项目，但它非常适合 coding agent：轻量、本地、Tree-sitter/WASM、SQLite、MCP，且工具设计围绕 agent 的上下文预算。OmniWeave 最初继承它，是正确路线。

\subsubsection{关键实现}

本地快照：\code{colbymchenry/codegraph@a893156}，见 \localpath{metadata/repo-snapshot.tsv}。

\begin{itemize}
  \item \textbf{工具 surface 极度收缩。} \code{src/mcp/tools.ts} 中 \code{getStaticTools} 默认只暴露 \code{codegraph\_explore}，其他工具通过环境 allowlist 才出现。这背后是非常强的 agent 产品判断：让模型少选错工具。
  \item \textbf{输出预算是核心逻辑。} \code{tools.ts} 有 inline output cap、adaptive budget、small repo 下减少 relationships/additional files、test file filtering 等策略。它把“别炸上下文”当成主逻辑，而不是响应末尾的截断。
  \item \textbf{解析热路径有安全阀。} \code{src/extraction/index.ts} 中有 parse timeout、worker thread、按文件数 recycle worker、batch IO、watch stale banner。WASM 内存无法 shrink 这类现实问题被纳入设计。
  \item \textbf{watcher 和 indexer scope 对齐。} \code{src/sync/watcher.ts} 对平台 watcher 复杂度、debounce、ignore matcher 复用、inotify/EMFILE degrade 都有处理。
  \item \textbf{callback synthesizer 走“高精度低召回”。} \code{src/resolution/callback-synthesizer.ts} 的动态边合成带 fan-out cap 和 provenance，而不是任意联想。
\end{itemize}

\subsubsection{OmniWeave 可借鉴}

codegraph 已经证明：\textbf{agent-first 的图不是“查询能力越多越好”，而是“一个 primary query 能不能把探索压成一跳”。} OmniWeave 应继续强化 \code{explore}/\code{node}，而不是把所有边种都暴露成独立工具。

建议：

\begin{enumerate}
  \item 将默认工具面进一步 A/B：是否可以默认只给 \code{explore}、\code{node}、\code{search}，把 \code{callers}/\code{impact} 作为 \code{explore} 内部路线或 allowlist。
  \item 把输出预算策略从“字符上限”推进到“任务形状预算”：single-point、reverse/blast、flow survey、file read replacement 各自有不同输出模板。
  \item 对所有新 edge kind 先问：它是否能显著减少 agent 工具调用？如果不能，就不要进默认输出。
\end{enumerate}

\subsection{Codebase-Memory：宏大、工程野心强，但需要防止过载}

\subsubsection{定位}

\project{Codebase-Memory} 是本轮最接近“代码库记忆后端”的系统：Tree-sitter 多语言、LSP hybrid、SQLite、FTS、Cypher-like 查询、MCP 14 tools、artifact 持久化、安装到多个 agent 的控制面。它的产品叙事很强：让 agent 不再反复重读整个 repo。

本地快照：\code{DeusData/codebase-memory-mcp@684e35b}。

\subsubsection{关键实现}

\begin{itemize}
  \item \textbf{14 个 MCP 工具。} \code{src/mcp/mcp.c:280-448} 定义 \code{index\_repository}、\code{search\_graph}、\code{query\_graph}、\code{trace\_path}、\code{get\_architecture}、\code{ingest\_traces}、ADR 等工具。这体现了“平台化”路线。
  \item \textbf{SQLite authorizer 值得重视。} \code{src/store/store.c:540-581} 注册 SQLite authorizer，拒绝 \code{ATTACH}/\code{DETACH} 等危险操作。对于任何暴露 query language 的图服务，这是底线。
  \item \textbf{RAM-first pipeline。} \code{src/pipeline/pipeline.c:1-12,621-754} 把 discover/structure/source/LZ4/definition/resolve/post-pass/dump 串成多阶段流程，先在内存图里做 pass，再 dump 到 SQLite。
  \item \textbf{自定义 SQLite 函数。} \code{store.c} 注册 \code{REGEXP}/\code{IREGEXP}、camel-case splitter、int8 cosine，FTS5 采用 contentless 表。这是把查询性能推入存储层的做法。
  \item \textbf{语义层很激进。} \code{src/semantic/semantic.c} 使用 TF-IDF/RRI/API/type/decorator/dataflow/graph diffusion 等信号，不一定读取源码，也不等同于 embedding API。这条路线可以做概念召回，但必须避免变成结构真相。
  \item \textbf{watcher 是 git-based。} \code{src/watcher/watcher.c} 走 \code{git status}/HEAD tracking，而不是 fsnotify；优点是跨平台和版本语义，缺点是对未纳入 git 的临时文件语义更弱。
\end{itemize}

\subsubsection{OmniWeave 可借鉴}

\begin{itemize}
  \item \textbf{artifact 共享。} Codebase-Memory 的 \code{.codebase-memory/graph.db.zst} 思路值得借鉴：OmniWeave 可以有 schema-versioned、hash-verified、read-only importable snapshot，用于 CI/团队冷启动。
  \item \textbf{安全查询层。} 如果 OmniWeave 未来支持 SQL/Cypher-like 查询，必须先有 authorizer、只读 AST、复杂度限制、row cap、timeout。
  \item \textbf{安装/控制面。} Codebase-Memory 和 code-review-graph 都在自动安装 MCP/skills/instructions 上用力。OmniWeave 的 installer 可以吸收“多 agent 目标 + dry-run + sibling config preservation”。
\end{itemize}

\subsubsection{不应照搬}

\begin{itemize}
  \item 14 tools 不适合作为 OmniWeave 默认形态。工具多会把 agent 重新带回“试错导航”。
  \item \code{query\_graph} 类自由查询若没有安全和输出治理，会把复杂度转嫁给 agent。
  \item \code{ingest\_traces} 这类运行时边如果没有验证闭环，容易制造“看起来很强、实际上不可维护”的边。
\end{itemize}

\subsection{CodeGraphContext/CGC：功能广，但复杂度已经明显外溢}

\subsubsection{定位}

\project{CGC} 把 code graph 做成多后端、多工具、多语言、bundle、watch、SCIP 可选、Spring/ORM/MyBatis 等 framework pass 的平台。它有很强的“企业代码理解平台”企图。

本地快照：\code{CodeGraphContext/CodeGraphContext@c235df7}。

\subsubsection{关键实现}

\begin{itemize}
  \item \textbf{工具面很大。} \code{tool\_definitions.py} 包含 indexing、find、relationship analysis、watch、Cypher、package graph、dead code、complexity、visualize、bundle、context switch、Java Spring endpoint/bean/datasource 等。
  \item \textbf{async indexing pipeline。} \code{tools/indexing/pipeline.py} 先 discover/pre-scan imports，再 semaphore 并发 parse，最后按确定顺序写图；之后再做 inheritance、CALLS、Spring、Maven/Gradle、ORM、MyBatis、embeddings 等 post-pass。
  \item \textbf{SCIP optional path。} \code{scip\_indexer.py} 能本地或 Docker 跑 SCIP indexer，针对 C/C++ 等依赖 compile commands 的语言 fallback 到 Tree-sitter。
  \item \textbf{ROADMAP 暴露出真实问题。} tool definitions 与 handler drift、Kuzu Cypher translation hack、bundle hash/signature 不足、graph builder recursion、16 language query toolkit stub 等，说明平台化会快速制造维护面。
\end{itemize}

\subsubsection{OmniWeave 可借鉴}

\begin{itemize}
  \item \textbf{可选 SCIP adapter。} 对 TypeScript/Java/Rust/C++ 等有成熟 SCIP indexer 的语言，OmniWeave 可以 ingest SCIP 作为 \code{provenance=scip} 的补充事实，不应把 SCIP 作为必须依赖。
  \item \textbf{bundle/context switch。} 大 monorepo 下允许 agent 显式切换 project graph 或 import bundle，有价值；但必须有 hash、schema version、签名/路径 guard。
  \item \textbf{framework pass 要证据门禁。} Spring/ORM/MyBatis 类 pass 可以借鉴，但只有在 false-positive gate 和 eval fixture 存在时才进核心。
\end{itemize}

\subsubsection{不应照搬}

CGC 的最大教训是：\textbf{多数据库后端、多 UI、多 framework、多 query language 不是天然高级，而是维护债。} OmniWeave 当前的 SQLite + MCP + CLI 主线更清晰。除非有明确用户收益和 eval，不能引入后端多态。

\subsection{Serena：LSP/IDE 能力层，不是持久图竞争者}

\subsubsection{定位}

\project{Serena} 是基于 LSP/JetBrains 的 coding-agent 工具层。它擅长 symbol overview、find symbol、find references、implementations、diagnostics、rename/safe delete/insert/replace 等 IDE 操作，也有 project memory。它不是持久图数据库，强项是“精确编辑能力”。

本地快照：\code{oraios/serena@dd7eb6d}。

\subsubsection{关键实现}

\begin{itemize}
  \item \code{agent.py} 里有 ToolSet/AvailableTools、include/exclude/fixed tool set、legacy name mapping、exclude editing tools 等工具模式控制。
  \item \code{ls\_manager.py} 并行启动 language servers，按文件选择合适 LS，失败可 restart。
  \item \code{symbol\_tools.py} 中 \code{GetSymbolsOverviewTool}、\code{FindSymbolTool}、\code{FindReferencingSymbolsTool} 都强调 compact result、body/depth/include filters、max answer chars。
  \item editing tools 以 symbol body 为单位替换、插入 before/after、rename、safe delete，路线更接近 IDE refactor。
\end{itemize}

\subsubsection{OmniWeave 可借鉴}

Serena 最值得借鉴的是 \textbf{tool mode}：read-only、editing、initial instructions、project activation、memory 各自可开关。OmniWeave 未来如果进入编辑/重构闭环，应把 LSP/Serena 类能力作为 companion，而不是让 graph core 承担 rename correctness。

建议：

\begin{enumerate}
  \item 增加 \code{validation mode}：OmniWeave 给出结构候选，LSP/SCIP 验证 same-language references。
  \item 对 MCP tools 提供 \code{read-only strict mode} 与 \code{edit-assist mode}，避免 agent 在纯研究时看到编辑工具。
  \item 借鉴 Serena 的 shortened output 与 ambiguity handling，但输出要保持 OmniWeave 的 source + blast radius 形态。
\end{enumerate}

\subsection{Aider repo-map：排名/压缩强，但没有可遍历关系}

\subsubsection{定位}

\project{Aider} 的 repo-map 是一个上下文压缩/排序系统，不是图数据库。它把 tree-sitter tags、defs/refs、聊天上下文、文件名/identifier mention、PageRank、token budget binary search 组合起来，产生“该给模型看的 repo map”。

本地快照：\code{Aider-AI/aider@5dc9490}。

\subsubsection{关键实现}

\begin{itemize}
  \item \code{aider/repomap.py} 缓存 tags，使用 tree-sitter query 文件和 Pygments fallback。
  \item \code{get\_ranked\_tags} 构建 \code{networkx.MultiDiGraph}，defs/refs 构边，按聊天文件、提到的文件名、identifier mention 个性化 PageRank。
  \item identifier 权重会提高更具体的 snake/kebab/camel 名称，降低 private/frequent identifiers。
  \item 输出通过 token budget binary search 调整到 \code{max\_map\_tokens}。
\end{itemize}

\subsubsection{OmniWeave 可借鉴}

OmniWeave 不应复制 repo-map 作为替代，但可以学习两个点：

\begin{enumerate}
  \item \textbf{Personalized ranking。} \code{explore} 可以把用户 query、当前文件、最近打开符号作为 personalization seed，对候选 symbol/source 排序。
  \item \textbf{Token-budget binary search。} 在输出临界区，用二分寻找最大可包含完整函数/完整文件段，而不是粗暴截断。
\end{enumerate}

边界：repo-map 没有 \code{callers(x)}、\code{impact(x)}、跨进程边、workflow artifact DAG，因此不是 OmniWeave 的核心竞品，而是输出排序参考。

\subsection{Codanna：Rust/Tantivy 符号索引，agent 消歧协议值得学}

\subsubsection{定位}

\project{Codanna} 是 Rust 写的 symbol index + relationship graph + semantic sidecar。它的工程重点在高性能 pipeline、Tantivy 搜索、\code{symbol\_id:N} 直接查找、MCP/CLI 双形态、watch/server lock。

本地快照：\code{bartolli/codanna@cb2c095}，本次使用 \localpath{repos/codanna-clean}；\localpath{repos/codanna} 是中断 checkout，不作为证据。

\subsubsection{关键实现}

\begin{itemize}
  \item \textbf{五阶段 pipeline。} \code{src/indexing/pipeline/mod.rs:1-20,260-543} 明确 DISCOVER → READ → PARSE → COLLECT → INDEX。READ/PARSE 多线程，COLLECT 单线程分配 ID，INDEX 单线程写 Tantivy，并用 bounded channels 做 backpressure。
  \item \textbf{Tantivy schema 分开 exact 与 ngram。} \code{src/storage/tantivy.rs:118-153} 同时有 \code{name} exact STRING 和 \code{name\_text} ngram TEXT，避免 \code{MyService} 匹配 \code{Main} 这类问题。
  \item \textbf{关系 metadata。} \code{relation\_receiver}、\code{relation\_static\_call} 等字段说明它在 call resolution 上处理 receiver/static call。
  \item \textbf{symbol\_id 消歧。} \code{src/retrieve.rs:66-86,119-190} 对 name lookup 的 ambiguous case 会返回 \code{symbol\_id:N} 建议；\code{src/mcp/mod.rs:84-100,299-322} 也把 symbol id 作为 MCP 参数。
  \item \textbf{向量侧车。} \code{src/vector/storage.rs} 用 mmap 二进制向量文件，自称 \textless 1us access；但 \code{read\_vector} 在片段中线性扫描 vector id，所以该指标需要结合 segment/index 设计验证，不能按注释直接采信。
\end{itemize}

\subsubsection{OmniWeave 可借鉴}

\begin{itemize}
  \item \textbf{直接 ID 工作流。} OmniWeave 的 \code{node}/\code{explore} 输出可以稳定暴露短 ID 或 continuation key，让后续工具无需重新按名字消歧。
  \item \textbf{exact + fuzzy 双字段。} SQLite FTS 已有基础，但 symbol name 的 exact/fuzzy/substring 语义应明确分层，避免误召回污染 agent。
  \item \textbf{pipeline telemetry。} Codanna 对 stage wait time、wall time、thread count 的度量值得吸收，用于 OmniWeave 的 indexing profile。
\end{itemize}

\subsection{Claude Context 与 semantic-search-mcp：语义入口，不是结构真相}

\subsubsection{Claude Context}

\project{Claude Context} 是 Zilliz 的 MCP semantic code search。README 明确需要 OpenAI embedding key 和 Milvus/Zilliz Cloud，目标是让 Claude Code 等 agent 从百万行代码中用语义搜索拿相关片段。它是“外部向量数据库 + embedding + MCP”的产品化路线。

本地快照：\code{zilliztech/claude-context@627eb2b}。

关键实现：

\begin{itemize}
  \item \code{packages/core/src/context.ts:131-184} 初始化 OpenAI embedding、vector database、AST splitter，默认支持一组扩展和 ignore patterns。
  \item \code{Context.getCollectionName} 用 codebase path hash 区分 collection，并 sanitize override，避免多个项目碰撞。
  \item \code{packages/mcp/src/handlers.ts:10-23} 用 \code{indexingTasks} 跟踪后台 indexing，clear index 时可以 cancel/await，避免“清空后后台继续写入”的竞态。
  \item \code{handlers.ts:31-154} 对 snapshot 0/0 completed poisoning 做启动修复，避免错误状态导致无限 reindex。
  \item \code{packages/mcp/src/sync.ts:59-180} 有全局 sync lock、stale lock recovery、background sync interval。
\end{itemize}

这些细节说明 Claude Context 的难点不是“embedding 一下代码”，而是外部向量库状态、snapshot、并发 index、同步锁、路径归一化、collection recovery。若 OmniWeave 引入语义 sidecar，必须把这些状态复杂度隔离在 optional 模块。

\subsubsection{semantic-search-mcp}

\project{semantic-search-mcp} 是更小的语义搜索基线：FastEmbed/Jina code embeddings、本地 SQLite、vector + FTS5 RRF、watcher、pause/resume/reindex/clear。它证明语义检索可以轻量，但仍只回答“哪里可能相关”，不回答“谁调用谁”。

\subsubsection{OmniWeave 可借鉴}

\begin{enumerate}
  \item 建一个 optional \code{semantic\_seed} 工具或 internal stage：自然语言 → candidate symbols/files → 交给 \code{explore} 做结构化上下文。
  \item 语义结果必须带 \code{source=semantic} 和低置信标签，不能被写入 \code{calls}/\code{crossLang} 等结构边。
  \item 不默认依赖外部 API 或 cloud vector DB；如果接入，必须是 plugin/sidecar，不污染 local-first 核心。
\end{enumerate}

\subsection{Blarify：LSP/SCIP + 图数据库 + LLM 文档生成，适合做反面边界}

\subsubsection{定位}

\project{Blarify} 目标是把本地 repo 表示成图，让 LLM traverses graph 理解逻辑；支持 Tree-sitter hierarchy、LSP/SCIP reference resolver、Neo4j/FalkorDB 管理、workflow/documentation 生成、MCP tools。

本地快照：\code{blarApp/blarify@ada9bbd}。

\subsubsection{关键实现}

\begin{itemize}
  \item \code{project\_graph\_creator.py:69-72} 先 create code hierarchy，再 create relationships from references。
  \item \code{project\_graph\_creator.py:170-246} 批量对所有 definition nodes 查 references，再用 relationship creator 建边。这是 LSP/SCIP 精确引用路线。
  \item \code{hybrid\_resolver.py:20-35,71-185} 提供 \code{SCIP\_ONLY}/\code{LSP\_ONLY}/\code{SCIP\_WITH\_LSP\_FALLBACK}/\code{AUTO}，Python/TypeScript 走 SCIP，否则 LSP。
  \item \code{prebuilt/graph\_builder.py:64-136} 把 graph build、DB save、workflow discovery、documentation generation、embedding 绑在一个 builder 上。
\end{itemize}

\subsubsection{OmniWeave 可借鉴与应避免}

可借鉴：SCIP/LSP 作为 reference resolution adapter；workflow discovery 作为后处理 pass 的思路；GraphEnvironment 这类环境 metadata。

应避免：把 LLM documentation、workflow generation、embedding、graph DB persistence 都塞进核心 builder。OmniWeave 的核心必须保持“静态事实 → 可信边 → agent 输出”，其余都应是插件。

\subsection{code-review-graph：垂直场景里的好产品化}

\subsubsection{定位}

\project{code-review-graph} 专注 PR/代码审查：Tree-sitter 结构图、SQLite、blast radius、风险评分、CI sticky comment、skills/slash commands。它不是通用 substrate，但在“review changes”场景里非常完整。

本地快照：\code{tirth8205/code-review-graph@b72413c}。

\subsubsection{关键实现}

\begin{itemize}
  \item \code{README.md:171-222} 很认真地区分 benchmark：528x 是最大值，不是典型值；recall 1.0 是 graph-derived upper bound，不是独立正确性。
  \item \code{docs/REPRODUCING.md:10-60} 校准 chars/4 token estimate，并提供 tiktoken verify；\code{docs/REPRODUCING.md:100-128} 给出完整 eval 命令和 failure semantics。
  \item \code{code\_review\_graph/graph.py:32-79} 用 SQLite 存 nodes/edges/metadata，edge 上有 \code{confidence}/\code{confidence\_tier}。
  \item \code{code\_review\_graph/tools/review.py:24-186} \code{get\_review\_context} 根据 changed files 算 impact radius、source snippets、review guidance 和 context savings。
  \item \code{parser.py:84-146} 支持大量扩展，包括 R、Perl、Lua、Objective-C、Bash、Jupyter、SQL、Vue/Svelte 等；但很多语言是浅结构抽取。
\end{itemize}

\subsubsection{OmniWeave 可借鉴}

\begin{itemize}
  \item \textbf{诚实 benchmark 写法。} 把最好情况、典型情况、circular upper bound、co-change independent-ish evidence 分开。这一点和 OmniWeave README 当前风格一致，应该延续。
  \item \textbf{review vertical package。} OmniWeave 可以提供 \code{review-delta} skill/command，但底层仍用 \code{impact}/\code{explore}，不要把 review 特化逻辑塞进 core graph。
  \item \textbf{context savings metadata。} 对每次 MCP 输出估算“替代 grep/read 节省了什么”，有助于 agent 和用户理解工具价值。
\end{itemize}

\subsection{code-graph-mcp 变体：同名项目很多，但参考价值有限}

本轮看了两个同名/近名项目：

\begin{itemize}
  \item \project{entrepeneur4lyf/code-graph-mcp@374d41e}：Python、ast-grep/rustworkx、多语言、usage guide、workflow orchestration。README 强调工具指导、性能等级、25+ 语言，但源码体量较小，更多是包装型工具。
  \item \project{sdsrss/code-graph-mcp@960d120}：Rust、Tree-sitter、SQLite/sqlite-vec、BM25+vector RRF、Merkle incremental、context compression、HTTP route tracing、Claude plugin。它的方向更接近“单 binary + hybrid search + route trace”。
\end{itemize}

可借鉴点：Rust 变体的 Merkle no-change detection、sqlite-vec optional embedding、route tracing、context compression level（L0/L1/L2/L3）值得单项评估。

不应照搬：README 中很多节省数字来自小型自测和项目自述，缺少 OmniWeave 这种 agent A/B 边界验证；不能作为战略证据。

\subsection{Graphify 与 Understand Anything：可视化学习层，而不是核心图引擎}

\project{Graphify} 与 \project{Understand Anything} 的共同方向是“把项目变成知识图/交互 dashboard/报告/导览”。它们对 onboarding、项目展示、非工程用户理解有价值，但不是低延迟 agent retrieval substrate。

Graphify 会生成 \code{graph.html}、\code{GRAPH\_REPORT.md}、\code{graph.json}，覆盖代码、docs、PDF、图片、视频，并提供各 agent 的 skill 安装。Understand Anything 用多 agent pipeline、interactive dashboard、domain view、guided tours、persona-adaptive UI。

OmniWeave 可以借鉴：

\begin{itemize}
  \item 独立导出 \code{architecture report}/\code{graph report}，作为人读产物。
  \item 给 onboarding 做可视化，不进入 MCP 默认路径。
  \item 把“surprising connections/suggested questions”作为 offline report，而不是实时工具噪声。
\end{itemize}

\subsection{SCIP/Sourcegraph：精确索引的事实格式，不是查询引擎本身}

\subsubsection{定位}

\project{SCIP} 是 Sourcegraph 推出的 code intelligence protocol。它不是数据库，也不是 agent 工具，而是 indexer 与 consumer 之间的 transmission format。

官方 README 说明它用于 Go to definition、Find references、Find implementations；\code{scip.proto} 说明 indexer 可处于 compiler precision 与 heuristic syntax analysis 之间。设计文档明确：SCIP 是 transmission format，不是 storage/query format，见 \localpath{repos/scip/docs/DESIGN.md:10-16,63-72}。

\subsubsection{关键设计}

\begin{itemize}
  \item \code{Index} 推荐 streaming consume，metadata 必须在 stream 开头，减少大 payload 内存压力，见 \localpath{repos/scip/scip.proto:20-39}。
  \item \code{Document.relative\_path} 规范要求相对 project root、不能 \code{/} 开头、不能符号链接、路径分隔符统一，见 \localpath{repos/scip/scip.proto:75-91}。
  \item Symbol string grammar 包含 scheme、package manager/name/version、descriptor，支持跨 repo/package 唯一性，见 \localpath{repos/scip/scip.proto:148-190}。
  \item 设计文档反对直接编码全图 adjacency list，理由是并行性、内存和 consumer 侧负担，见 \localpath{repos/scip/docs/DESIGN.md:97-123}。
\end{itemize}

\subsubsection{OmniWeave 可借鉴}

\begin{enumerate}
  \item \textbf{SCIP importer，而非 SCIP core。} OmniWeave 可读取 \code{index.scip}，导入 precise same-language references/inheritance，标注 \code{provenance=scip}。
  \item \textbf{Symbol ID 规范。} OmniWeave 当前 ID 偏本地 DB；未来 snapshot/bundle/export 可参考 SCIP 的 package-aware symbol scheme。
  \item \textbf{Streaming artifact。} 如果做团队共享索引，不要直接 dump 全图 JSON；应支持 streaming、version、root path、position encoding。
\end{enumerate}

\subsection{Kythe/Glean/CodeQL/Joern/Semgrep/OpenGrep：重型工业系统的启示}

这些系统没有全部本地 clone 深读，本报告基于官方文档/公开资料定位：

\begin{itemize}
  \item \project{Kythe}：强调跨语言 semantic code graph、facts/edges、serving tables、cross-reference。适合大规模公司代码导航，不适合轻量 agent MCP core。
  \item \project{Glean}：Meta 开源的代码索引系统，使用 Angle query language，目标是多语言代码索引与查询。它代表“Datalog-like facts + query”的工业路线。
  \item \project{CodeQL}：GitHub 的 semantic code analysis engine，把代码作为数据库查询，主要用于 variant analysis/security。强在安全分析，不在低延迟 agent 上下文。
  \item \project{Joern}：Code Property Graph，AST/CFG/PDG 等统一为 CPG，用 Scala DSL 挖漏洞模式。适合安全研究，不适合作为日常 agent exploration 默认层。
  \item \project{Semgrep/OpenGrep}：规则/模式匹配、SAST、taint/reachability。适合“找安全/规范模式”，不是 repo understanding graph。
\end{itemize}

结论：OmniWeave 应借鉴它们的\textbf{事实模型纪律、query safety、provenance、版本化 schema、可解释规则}，而不是照搬重型查询语言或安全扫描平台。

\subsection{RepoGraph 与 CodeRAG-Bench：研究评测提醒我们别只测检索}

\project{RepoGraph} 把 repo-level code graph 加入 SWE-bench/Agentless/SWE-agent。源码显示它大量借鉴 Aider repo-map、Agentless、grep-ast；\code{construct\_graph.py} 中甚至有对 Python 代码做 \code{exec(import_statement)} 来识别标准库函数的逻辑，风险很高，不适合作为生产参考。

\project{CodeRAG-Bench} 关注“检索增强是否提升代码生成”。它的 retrieval 脚本覆盖 BM25、SBERT、OpenAI/Voyage embeddings、repoeval、swe-bench repo retrieval，generation 侧再评估模型输出。这对 OmniWeave 的启发不是实现，而是评测：只测 Recall@K 或 token saving 不够，还要测最终 agent task 是否更快、更稳、更少误读。

\section{横向比较}

\subsection{事实来源：Tree-sitter、LSP、SCIP、embedding 的信任边界}

\begin{longtable}{p{0.18\linewidth}p{0.28\linewidth}p{0.23\linewidth}p{0.23\linewidth}}
\toprule
来源 & 强项 & 弱点 & OmniWeave 应用 \\
\midrule
\endhead
Tree-sitter/AST & 零配置、快、跨语言、可本地运行 & 类型/动态分派/宏/导入解析有限 & 核心 extractor，配 provenance/confidence。 \\
LSP & 同语言精确 references/rename/diagnostics & 需要环境、慢启动、跨语言弱 & optional validation/edit companion。 \\
SCIP & 编译器/indexer 级别 precise facts，可跨 repo symbol & 依赖 indexer、生成成本高，不是查询引擎 & importer/snapshot fact source。 \\
Embedding/vector & 概念召回强，自然语言入口 & 不产生结构真相，可解释弱，需状态管理 & optional semantic seed。 \\
规则/CPG/CodeQL & 安全/variant analysis 强 & 规则语言重、学习成本高 | 安全插件/专项 pass，不进默认路径。 \\
\bottomrule
\end{longtable}

\subsection{存储与查询}

\begin{itemize}
  \item \textbf{SQLite 是 OmniWeave 当前最优核心。} codegraph、OmniWeave、code-review-graph、Codebase-Memory、sdsrss/code-graph-mcp 都选择 SQLite 或 SQLite 衍生能力。它足够本地、可携带、低运维。
  \item \textbf{Tantivy/FTS/vector 是 sidecar 候选。} Codanna 的 Tantivy exact/ngram 值得学；sqlite-vec/FastEmbed 适合 optional semantic layer。
  \item \textbf{Neo4j/Falkor/Kuzu 多后端不适合核心。} CGC/Blarify 的多后端路线会带来 schema parity、query translation、deployment、错误处理复杂度。
  \item \textbf{自由 query language 必须先安全化。} Codebase-Memory 的 SQLite authorizer 是底线；没有 authorizer、timeout、row cap、explain/complexity gate，就不应给 agent 自由查询。
\end{itemize}

\subsection{增量同步}

生态里主要有五种：

\begin{enumerate}
  \item \textbf{文件 watcher + debounce}：codegraph/OmniWeave、semantic-search-mcp；适合开发态，但要处理平台 watcher 差异和 stale state。
  \item \textbf{Git diff/status}：Codebase-Memory、code-review-graph；适合 review/CI，但对未跟踪文件不完整。
  \item \textbf{Merkle/hash tree}：sdsrss/code-graph-mcp；no-change detection 快，适合大 repo。
  \item \textbf{snapshot artifact}：Codebase-Memory；适合团队共享和冷启动。
  \item \textbf{remote/cloud collection sync}：Claude Context；状态复杂度最高，必须有 lock/recovery。
\end{enumerate}

OmniWeave 当前 watcher + staleness banner 是合理核心；下一步最值得加的是 \textbf{hash-verified snapshot import/export}，不是云同步。

\subsection{MCP 工具面}

\begin{longtable}{p{0.22\linewidth}p{0.34\linewidth}p{0.34\linewidth}}
\toprule
项目 & 工具面趋势 & 对 agent 的风险 \\
\midrule
\endhead
codegraph & 默认只暴露 primary \code{explore}，其他 opt-in & 低，模型选择成本小。 \\
OmniWeave & 默认五个核心工具，强调 \code{explore}/\code{node} & 可继续收缩或按 query shape routing。 \\
Codebase-Memory & 14 tools，含 query/architecture/ADR/runtime traces & 能力强但选择成本高。 \\
CGC & indexing、watch、Cypher、visualize、bundle、framework tools & 极易 tool drift 和输出面不一致。 \\
Serena & tool set/mode 可控，读写工具分层 & 值得借鉴 mode control。 \\
code-review-graph & review/change vertical tools 很完整 & 场景清晰时有效，通用时偏重。 \\
\bottomrule
\end{longtable}

结论：OmniWeave 的主线应是 \textbf{少工具、强输出、明确 fallback}。每增加一个 tool，必须证明它减少了 agent 的实际工具调用或错误选择。

\section{OmniWeave 的优势：不是“大而全”，而是边界可信}

\subsection{优势一：抓住 LSP/grep 共同盲区}

grep 的盲区是“关系”；LSP 的盲区是“跨语言/跨进程/工作流/动态框架约定”。大部分竞品只解决其中之一：

\begin{itemize}
  \item Aider 解决“哪些文件值得读”，不是“谁调用谁”。
  \item Serena/LSP 解决同语言 symbol navigation，不解决 Python→R、Snakemake→tool。
  \item Claude Context 解决“概念上相关”，不解决 structure。
  \item CodeQL/Joern 解决安全程序分析，不解决 agent 日常上下文形态。
\end{itemize}

OmniWeave 的 \code{crossLang}、S4 dispatch、workflow DAG、external tool nodes 正好打在这些交集盲区。这个定位应坚持。

\subsection{优势二：诚实边界比高召回更稀缺}

很多项目把“能猜到”当成“能连边”。OmniWeave 当前更克制：

\begin{itemize}
  \item 动态路径里 residual \code{\{...\}}/\code{\$} 就 skip；
  \item interpreter 必须处在真实 subprocess call 里，不能只是 echo 字符串；
  \item fan-out 有 cap；
  \item heuristic edge 带 confidence；
  \item static dispatch 只建声明骨架，不伪造运行时 target。
\end{itemize}

这正是 agent substrate 的信任基础。Agent 会把结构边当成事实来行动；错边会导致错误修改、错误影响面、错误审查结论。OmniWeave 的“宁可漏，不可骗”是优势，不是保守。

\subsection{优势三：A/B 证明了真实价值边界}

README 当前最有力量的不是数字本身，而是承认边界：correctness 经常 tie；single-point lookup 不赢；concept search 不是 OmniWeave；cross-process at scale 由于运行时路径构造也会 tie。很多竞品还停在“我减少了多少 token”的自述阶段，OmniWeave 已经进入“在哪些 query shape 赢、在哪些不赢”的阶段。

这会带来一个战略好处：产品路线可以少走弯路。比如：

\begin{itemize}
  \item 不应该把资源投向“让 symbol lookup 比 grep 更快很多”；
  \item 应该投向 reverse/blast-radius、workflow/crossLang、weak-model guardrail、output sufficiency；
  \item 不应该把 semantic search 写成核心卖点；
  \item 可以把 semantic search 作为 seed layer 服务 \code{explore}。
\end{itemize}

\subsection{优势四：本地、轻量、可分发}

Codebase-Memory、Claude Context、CGC、Blarify 都有更重的状态面：外部 vector DB、多 DB backend、SCIP/Docker/LSP、cloud collection、Neo4j/Falkor。OmniWeave 的核心依赖更少：Node 22.5+、WASM tree-sitter、SQLite。这是安装、CI、隐私、团队采用的实际优势。

但要注意：轻量不是简陋。轻量必须和 artifact、version、integrity、watch freshness、schema migration 一起做，否则只是在本地临时可用。

\section{OmniWeave 应该借鉴什么}

\subsection{优先级 P0：输出面与工具路由继续压实}

\begin{longtable}{p{0.2\linewidth}p{0.32\linewidth}p{0.38\linewidth}}
\toprule
借鉴对象 & 做法 & OmniWeave 落地建议 \\
\midrule
\endhead
codegraph & 默认只暴露 primary tool，输出预算按 repo/task 调整 & 对默认 tools 做 A/B：\code{explore}+\code{node}+\code{search} 是否足够；其余 opt-in。 \\
Aider & PageRank + mention personalization & \code{explore} 候选排序引入 query/current-file/recent-symbol personalization。 \\
code-review-graph & context savings 元数据与诚实 benchmark 叙事 & MCP 响应可附 compact savings/debug metadata，仅在 verbose 或 eval mode 展示。 \\
Codanna & \code{symbol\_id:N} 消歧协议 & \code{node}/\code{explore} 输出 stable continuation key，后续工具直接消费。 \\
\bottomrule
\end{longtable}

\subsection{优先级 P1：snapshot/artifact 共享}

借鉴 Codebase-Memory 的 artifact 方向，但做得更克制：

\begin{itemize}
  \item \code{omniweave snapshot export}：输出 \code{graph.db.zst} 或 bundle directory；
  \item manifest 包含 schema version、OmniWeave version、source root hash、file count、edge kind count、createdAt、platform；
  \item import 默认 read-only，路径 root 必须显式 remap；
  \item 不信任 artifact 内任意路径；所有 source snippet 读取仍受 project root guard；
  \item snapshot import 只加速冷启动，不替代本地 reindex 的真相。
\end{itemize}

\subsection{优先级 P1：可选 SCIP importer}

SCIP 是最合理的 precise same-language facts 来源。建议：

\begin{enumerate}
  \item 首期只支持读取 \code{index.scip}，不负责运行 indexer。
  \item 只导入 definition/reference/implementation/inheritance 等清晰事实，标注 \code{provenance=scip}。
  \item 与 Tree-sitter 边冲突时不覆盖，保留 parallel provenance 或 higher confidence。
  \item 输出中对 SCIP facts 标注 compiler-backed/third-party indexed。
\end{enumerate}

这样可以补 OmniWeave 在 same-language 精度上的短板，但不会让安装路径变重。

\subsection{优先级 P2：semantic sidecar}

借鉴 Claude Context、semantic-search-mcp、Cursor 的方向，但不做默认依赖：

\begin{itemize}
  \item semantic sidecar 只返回 candidate file/symbol IDs；
  \item \code{explore} 用这些 IDs 进入结构图；
  \item 语义层可以支持 local embedding、sqlite-vec 或外部 provider，但均为 plugin；
  \item 任何 semantic score 不能变成 structural edge confidence。
\end{itemize}

\subsection{优先级 P2：review vertical skill}

借鉴 code-review-graph 的场景化：

\begin{itemize}
  \item \code{/omniweave:review-delta}：自动取 git diff changed files → impact → affected workflows/tools → test gaps；
  \item \code{/omniweave:review-pr}：base branch diff → risk summary；
  \item 输出必须短、稳定、可粘贴，不暴露整个 graph JSON；
  \item 所有 review 指标都要标明 graph-derived 上界还是 git co-change evidence。
\end{itemize}

\section{OmniWeave 不应该借鉴什么}

\subsection{不要默认追求大工具面}

Codebase-Memory 的 14 tools、CGC 的大工具面、code-review-graph 的 vertical tools 都各有道理，但 OmniWeave 的默认 MCP surface 必须克制。对 agent 来说，多一个工具不是多一个能力，而是多一个误选路径。

判断标准：

\begin{quote}
  如果一个新增 tool 不能在 A/B 中降低工具调用、降低 token、减少 fallback read、或提升稳定性，就不应默认暴露。
\end{quote}

\subsection{不要引入多数据库后端}

Neo4j、FalkorDB、Kuzu、Ladybug、remote graph DB 都可以做 enterprise story，但会破坏 OmniWeave 的 local-first 和零配置优势。SQLite 已足够支持当前核心；更值得做的是 schema、index、query plan、snapshot，而不是 backend abstraction。

\subsection{不要把 LLM 文档生成塞进核心}

Blarify、Graphify、Understand Anything 都做 documentation/report/tour。这些适合 offline product layer，不适合 indexer core。LLM 生成内容不是事实，必须和 graph facts 分开存储、分开 provenance。

\subsection{不要把语义搜索当结构事实}

向量检索可以帮助发现入口，但它不能回答调用关系、影响半径、workflow artifact、S4 dispatch。语义层若进入核心边表，会破坏 graph trust。

\subsection{不要复制安全分析平台}

CodeQL/Joern/Semgrep/OpenGrep 很强，但它们解决的是 variant analysis、CPG mining、taint/rule matching。OmniWeave 可在未来以 plugin 接入安全 pass，但不应把 query DSL 或 taint engine 变成 MCP 默认能力。

\section{建议路线图}

\subsection{Phase A：输出与工具面硬化}

\begin{enumerate}
  \item 为 \code{explore}/\code{node}/\code{callers}/\code{impact} 建 query-shape goldens：single-point、reverse、flow survey、file read replacement、crossLang、workflow。
  \item 每个工具输出记录：字符数、source snippets 完整率、omitted reason、fallback read rate。
  \item A/B 默认工具面：5 tools vs 3 tools vs explore-only，指标为工具调用、token、turns、正确率、模型误选率。
  \item 引入 continuation key：一个 symbol/file/flow 被返回后，后续工具可用 key 直接定位。
\end{enumerate}

Done when：默认工具面比当前不增加 correctness 风险，同时降低平均 tool mis-pick 或 fallback read。

\subsection{Phase B：snapshot artifact}

\begin{enumerate}
  \item 定义 \code{omniweave-snapshot.json} manifest。
  \item 支持 export/import read-only snapshot。
  \item 引入 schema version and compatibility check。
  \item 做 corrupt snapshot、wrong root、stale source、malicious path regression tests。
\end{enumerate}

Done when：干净 checkout 可在无全量 reindex 的情况下读 snapshot，但任何源码变更都能触发 stale warning/reindex path。

\subsection{Phase C：SCIP importer}

\begin{enumerate}
  \item 读取 SCIP protobuf，导入 definitions/references。
  \item 映射到 OmniWeave nodes/edges，保留 original symbol string。
  \item 输出中展示 \code{provenance=scip}。
  \item 与 Tree-sitter edge 冲突时不静默覆盖。
\end{enumerate}

Done when：在 TS/Java/Rust fixture 上 same-language references 明显更精确，且不影响 crossLang/workflow/S4 gates。

\subsection{Phase D：semantic sidecar}

\begin{enumerate}
  \item 先做 internal seed，不暴露为默认 MCP tool。
  \item 本地 provider 优先，外部 provider 明确 opt-in。
  \item Eval 只测“是否帮助 concept query 找到正确 explore seed”，不把它混入 structural correctness。
\end{enumerate}

Done when：concept query 的 first useful file rate 提升，且 structure query 无 regression。

\subsection{Phase E：review vertical skill}

\begin{enumerate}
  \item 用现有 \code{impact}/workflow/tool edges 生成 PR review context。
  \item 输出 risk/test gaps/affected tools/affected workflows。
  \item 区分 graph-derived upper bound 和 git co-change evidence。
  \item 不默认安装，作为 skill/plugin 分发。
\end{enumerate}

\section{最终判断}

OmniWeave 应该避免两个诱惑：一个是向上变成“通用知识图/可视化平台”，一个是向下变成“更快 grep/更多语言 parser”。前者会拖垮核心，后者会失去差异化。

最强路径是：

\begin{quote}
  \textbf{把 OmniWeave 做成 coding agent 的可信结构控制层：少工具、强输出、边界可信、跨语言/跨进程/workflow/dynamic-dispatch 优先、本地可分发，并用 A/B 证明每个新增能力真的减少 agent 成本。}
\end{quote}

这也是与所有竞品拉开距离的地方。Codebase-Memory 更像“代码库记忆平台”；Serena 更像“IDE/LSP agent 工具层”；Claude Context 更像“语义召回层”；CodeQL/Joern 是“安全分析层”；Graphify/Understand Anything 是“学习/可视化层”。OmniWeave 的最佳位置不是吞并它们，而是成为它们之下或旁边的可信结构骨架。

\appendix

\section{本地源码快照}

\begin{longtable}{p{0.2\linewidth}p{0.16\linewidth}p{0.12\linewidth}p{0.12\linewidth}p{0.3\linewidth}}
\toprule
项目 & commit & 文件数 & 体量(KB) & 备注 \\
\midrule
\endhead
codegraph & a893156 & 387 & 20076 & OmniWeave 直接基座。 \\
codebase-memory & 684e35b & 1676 & 1280540 & C/Tree-sitter/SQLite/MCP，含 vendored grammars。 \\
cgc & c235df7 & 998 & 116640 & 多后端、多工具、SCIP optional。 \\
serena & dd7eb6d & 906 & 12104 & LSP/IDE 工具层。 \\
aider & 5dc9490 & 691 & 159884 & repo-map/ranking/context compression。 \\
codanna-clean & cb2c095 & 822 & 15800 & Rust/Tantivy/vector/symbol-id。 \\
claude-context & 627eb2b & 177 & 35092 & semantic search + Milvus/Zilliz。 \\
blarify & ada9bbd & 400 & 14220 & Tree-sitter/LSP/SCIP/graph DB/LLM docs。 \\
code-review-graph & b72413c & 287 & 15520 & PR review vertical。 \\
code-graph-mcp & 374d41e & 30 & 868 & Python MCP graph variant。 \\
code-graph-mcp-2 & 960d120 & 224 & 5540 & Rust MCP graph variant。 \\
semantic-search-mcp & 40a1b01 & 48 & 976 & local semantic search baseline。 \\
scip & 75d1151 & 219 & 3728 & protocol repo。 \\
graphify & a4d09ae & 592 & 14144 & report/visual graph layer。 \\
understand-anything & 7f5a717 & 420 & 51304 & plugin/dashboard/learning layer。 \\
repograph & 6c3977d & 209 & 7416 & research graph for SWE-bench. \\
coderag-bench & f9e100ca & 27674 & 1244536 & retrieval/generation benchmark. \\
\bottomrule
\end{longtable}

\section{关键证据索引}

\begin{longtable}{p{0.26\linewidth}p{0.64\linewidth}}
\toprule
证据 & 用途 \\
\midrule
\endhead
\localpath{OmniWeave/README.md:35-98} & OmniWeave agent A/B 边界、wins/ties/losses。 \\
\localpath{OmniWeave/README.md:135-164} & crossLang、S4、workflow DAG、external tool graph。 \\
\localpath{OmniWeave/CHECKPOINT.md:38-55} & 数据流与 skip-over-guess 边界规则。 \\
\localpath{OmniWeave/src/resolution/callback-synthesizer.ts:1941-2039} & general cross-language edge gates。 \\
\localpath{codegraph/src/mcp/tools.ts} & primary explore/output budget/default tool surface。 \\
\localpath{codegraph/src/extraction/index.ts} & parse worker、timeout、recycle、watch integration。 \\
\localpath{codebase-memory/src/mcp/mcp.c:280-448} & 14 MCP tools 与平台化工具面。 \\
\localpath{codebase-memory/src/store/store.c:540-581} & SQLite authorizer/security。 \\
\localpath{codebase-memory/src/pipeline/pipeline.c} & RAM-first multi-pass pipeline。 \\
\localpath{cgc/src/codegraphcontext/tools/indexing/pipeline.py} & async indexing + framework post-passes。 \\
\localpath{serena/src/serena/tools/symbol_tools.py} & LSP symbol/find/reference tools。 \\
\localpath{aider/aider/repomap.py} & PageRank/personalization/token budget repo-map。 \\
\localpath{codanna-clean/src/retrieve.rs:66-190} & \code{symbol\_id:N} 消歧工作流。 \\
\localpath{codanna-clean/src/indexing/pipeline/mod.rs:260-543} & DISCOVER/READ/PARSE/COLLECT/INDEX pipeline。 \\
\localpath{claude-context/packages/mcp/src/handlers.ts:10-154} & indexing task cancel/snapshot poisoning recovery。 \\
\localpath{blarify/blarify/code_references/hybrid_resolver.py:20-185} & SCIP/LSP hybrid resolver。 \\
\localpath{code-review-graph/docs/REPRODUCING.md:10-60,100-128} & token estimate calibration/eval reproducibility。 \\
\localpath{scip/scip.proto:20-39,75-190} & streaming index/document/symbol schema。 \\
\localpath{scip/docs/DESIGN.md:10-16,97-123} & SCIP 是 transmission format，避免全图 adjacency。 \\
\bottomrule
\end{longtable}

\section{外部资料与链接}

\begin{thebibliography}{99}
\bibitem{codegraph} colbymchenry/codegraph. \url{https://github.com/colbymchenry/codegraph}
\bibitem{codebase-memory} DeusData/codebase-memory-mcp. \url{https://github.com/DeusData/codebase-memory-mcp}
\bibitem{codebase-paper} Codebase-Memory paper, arXiv:2603.27277. \url{https://arxiv.org/abs/2603.27277}
\bibitem{serena} oraios/serena. \url{https://github.com/oraios/serena}
\bibitem{aider} Aider-AI/aider. \url{https://github.com/Aider-AI/aider}
\bibitem{codanna} bartolli/codanna. \url{https://github.com/bartolli/codanna}
\bibitem{claude-context} zilliztech/claude-context. \url{https://github.com/zilliztech/claude-context}
\bibitem{blarify} blarApp/blarify. \url{https://github.com/blarApp/blarify}
\bibitem{code-review-graph} tirth8205/code-review-graph. \url{https://github.com/tirth8205/code-review-graph}
\bibitem{scip} SCIP repository. \url{https://github.com/scip-code/scip}
\bibitem{scip-blog} Sourcegraph, SCIP announcement. \url{https://sourcegraph.com/blog/announcing-scip}
\bibitem{sourcegraph-ci} Sourcegraph code intelligence docs. \url{https://docs.sourcegraph.com/user/code_intelligence}
\bibitem{kythe} Kythe storage documentation. \url{https://kythe.io/docs/kythe-storage.html}
\bibitem{glean} Meta Engineering, Glean open source code indexing. \url{https://engineering.fb.com/2024/12/19/developer-tools/glean-open-source-code-indexing/}
\bibitem{codeql} GitHub CodeQL overview. \url{https://codeql.github.com/docs/codeql-overview/about-codeql/}
\bibitem{joern} Joern Code Property Graph docs. \url{https://docs.joern.io/code-property-graph/}
\bibitem{semgrep} Semgrep documentation. \url{https://docs.semgrep.dev/}
\bibitem{opengrep} OpenGrep repository. \url{https://github.com/opengrep/opengrep}
\bibitem{cursor-indexing} Cursor secure codebase indexing. \url{https://cursor.com/blog/secure-codebase-indexing}
\bibitem{cursor-search} Cursor semantic and agentic search docs. \url{https://cursor.com/docs/agent/tools/search}
\bibitem{github-copilot-indexing} GitHub Copilot repository indexing docs. \url{https://docs.github.com/en/copilot/concepts/context/repository-indexing}
\bibitem{openai-codex-skills} OpenAI Codex Agent Skills docs. \url{https://developers.openai.com/codex/skills}
\bibitem{openai-codex} OpenAI Codex repository. \url{https://github.com/openai/codex}
\bibitem{repograph} RepoGraph. \url{https://github.com/ozyyshr/RepoGraph}
\bibitem{coderagbench} CodeRAG-Bench. \url{https://github.com/code-rag-bench/code-rag-bench}
\end{thebibliography}

\end{document}
